% GENERATED FILE. Edit canonical lessons, then run npm run generate:books.
% canonical-source-hash: fnv1a64:8dce6b383c6641a1
% canonical-lessons: TA-C119-cavi, TA-C119-puttu, TA-C119-tuntu, TA-C119-talaiyanai, TA-C119-porvai

\chapter{Key, Lock, Towel, Pillow, Blanket}
\label{ch:ta-key-lock-towel-pillow-blanket}
\hlchaptermodality{119}

\begin{chapteropening}
\textbf{By the end of this chapter:} \emph{I can say a key, a lock, a towel, a pillow and a blanket.}

Complete the last lesson of chapter 119: i can say a key, a lock, a towel, a pillow and a blanket.
\end{chapteropening}

\section[\texorpdfstring{\ta{சாவி} (cāvi)}{சாவி (cāvi)}]{\ta{சாவி} (cāvi) — a key}

\label{lesson:TA-C119-cavi}



\begin{quote}
Before the new one: say the Tamil for a cup, then the Tamil for a spoon.
\end{quote}

\subsection*{You'll want to know: \ta{சாவி}}
\textbf{\ta{சாவி}} — \emph{cāvi} — \textquotedblleft{}a key\textquotedblright{}.

\textbf{In the family, and next door}

\noindent
\begin{tabularx}{\linewidth}{@{}>{\raggedright\arraybackslash}X>{\raggedright\arraybackslash}X@{}}
\toprule
\textbf{} & \textbf{word} \\
\midrule
Telugu & \te{తాళంచెవి} (\emph{tāḷaṁcevi}) \\
Malayalam & \ml{താക്കോൽ} (\emph{tākkōl}) \\
English & a key \\
\bottomrule
\end{tabularx}

\begin{grammarlens}[title={What you've built}]
One more word for this chapter.
\end{grammarlens}

\subsection*{Your turn}
\begin{itemize}
\raggedright
  \item \emph{Say it:} \emph{cāvi}
  \item \emph{Say it:} \emph{cāvi}, once more
  \item \emph{Say it:} say \emph{cāvi}
  \item \emph{Recall:} say the Tamil for a cup, then the Tamil for a spoon, then say \emph{cāvi} again
\end{itemize}

\begin{tcolorbox}[breakable,colback=teal!4,colframe=teal!35!black,title={Before you move on}]
What does \ta{சாவி} mean? (\textquotedblleft{}A key\textquotedblright{}.) Say it once more.
\end{tcolorbox}
\section[\texorpdfstring{\ta{பூட்டு} (pūṭṭu)}{பூட்டு (pūṭṭu)}]{\ta{பூட்டு} (pūṭṭu) — a lock}

\label{lesson:TA-C119-puttu}



\begin{quote}
Before the new one: say the Tamil for a spoon, then the Tamil for a key.
\end{quote}

\subsection*{You'll want to know: \ta{பூட்டு}}
\textbf{\ta{பூட்டு}} — \emph{pūṭṭu} — \textquotedblleft{}a lock\textquotedblright{}.

\textbf{In the family, and next door}

\noindent
\begin{tabularx}{\linewidth}{@{}>{\raggedright\arraybackslash}X>{\raggedright\arraybackslash}X@{}}
\toprule
\textbf{} & \textbf{word} \\
\midrule
Kannada & \ka{ಬೀಗ} (\emph{bīga}) \\
Telugu & \te{తాళం} (\emph{tāḷaṁ}) \\
Hindi (neighbour) & \dv{ताला} (\emph{tālā}) \\
English & a lock \\
\bottomrule
\end{tabularx}

\begin{grammarlens}[title={What you've built}]
One more word for this chapter.
\end{grammarlens}

\subsection*{Your turn}
\begin{itemize}
\raggedright
  \item \emph{Say it:} \emph{pūṭṭu}
  \item \emph{Say it:} \emph{pūṭṭu}, once more
  \item \emph{Say it:} say \emph{pūṭṭu}
  \item \emph{Recall:} say the Tamil for a spoon, then the Tamil for a key, then say \emph{pūṭṭu} again
\end{itemize}

\begin{tcolorbox}[breakable,colback=teal!4,colframe=teal!35!black,title={Before you move on}]
What does \ta{பூட்டு} mean? (\textquotedblleft{}A lock\textquotedblright{}.) Say it once more.
\end{tcolorbox}
\section[\texorpdfstring{\ta{துண்டு} (tuṇṭu)}{துண்டு (tuṇṭu)}]{\ta{துண்டு} (tuṇṭu) — a towel; a piece}

\label{lesson:TA-C119-tuntu}



\begin{quote}
Before the new one: say the Tamil for a key, then the Tamil for a lock.
\end{quote}

\subsection*{You'll want to know: \ta{துண்டு}}
\textbf{\ta{துண்டு}} — \emph{tuṇṭu} — \textquotedblleft{}a towel; a piece\textquotedblright{}.

\textbf{In the family, and next door}

\noindent
\begin{tabularx}{\linewidth}{@{}>{\raggedright\arraybackslash}X>{\raggedright\arraybackslash}X@{}}
\toprule
\textbf{} & \textbf{word} \\
\midrule
Kannada & \ka{ಟವೆಲ್} (\emph{ṭavel}) \\
Telugu & \te{తువ్వాలు} (\emph{tuvvālu}) \\
Malayalam & \ml{തോർത്ത്} (\emph{tōrttŭ}) \\
Hindi (neighbour) & \dv{तौलिया} (\emph{tauliyā}) \\
English & a towel \\
\bottomrule
\end{tabularx}

\begin{grammarlens}[title={What you've built}]
One more word for this chapter.
\end{grammarlens}

\subsection*{Your turn}
\begin{itemize}
\raggedright
  \item \emph{Say it:} \emph{tuṇṭu}
  \item \emph{Say it:} \emph{tuṇṭu}, once more
  \item \emph{Say it:} say \emph{tuṇṭu}
  \item \emph{Recall:} say the Tamil for a key, then the Tamil for a lock, then say \emph{tuṇṭu} again
\end{itemize}

\begin{tcolorbox}[breakable,colback=teal!4,colframe=teal!35!black,title={Before you move on}]
What does \ta{துண்டு} mean? (\textquotedblleft{}A towel; a piece\textquotedblright{}.) Say it once more.
\end{tcolorbox}
\section[\texorpdfstring{\ta{தலையணை} (talaiyaṇai)}{தலையணை (talaiyaṇai)}]{\ta{தலையணை} (talaiyaṇai) — a pillow}

\label{lesson:TA-C119-talaiyanai}



\begin{quote}
Before the new one: say the Tamil for a lock, then the Tamil for a towel; a piece.
\end{quote}

\subsection*{You'll want to know: \ta{தலையணை}}
\textbf{\ta{தலையணை}} — \emph{talaiyaṇai} — \textquotedblleft{}a pillow\textquotedblright{}.

\textbf{In the family, and next door}

\noindent
\begin{tabularx}{\linewidth}{@{}>{\raggedright\arraybackslash}X>{\raggedright\arraybackslash}X>{\raggedright\arraybackslash}X@{}}
\toprule
\textbf{} & \textbf{word} & \textbf{same root} \\
\midrule
Kannada & \ka{ದಿಂಬು} (\emph{dimbu}) &  \\
Telugu & \te{దిండు} (\emph{diṇḍu}) &  \\
Malayalam & \ml{തലയണ} (\emph{talayaṇa}) & yes \\
Hindi (neighbour) & \dv{तकिया} (\emph{takiyā}) &  \\
English & a pillow &  \\
\bottomrule
\end{tabularx}

\begin{grammarlens}[title={What you've built}]
One more word for this chapter.
\end{grammarlens}

\subsection*{Your turn}
\begin{itemize}
\raggedright
  \item \emph{Say it:} \emph{talaiyaṇai}
  \item \emph{Say it:} \emph{talaiyaṇai}, once more
  \item \emph{Say it:} say \emph{talaiyaṇai}
  \item \emph{Recall:} say the Tamil for a lock, then the Tamil for a towel; a piece, then say \emph{talaiyaṇai} again
\end{itemize}

\begin{tcolorbox}[breakable,colback=teal!4,colframe=teal!35!black,title={Before you move on}]
What does \ta{தலையணை} mean? (\textquotedblleft{}A pillow\textquotedblright{}.) Say it once more.
\end{tcolorbox}
\section[\texorpdfstring{\ta{போர்வை} (pōrvai)}{போர்வை (pōrvai)}]{\ta{போர்வை} (pōrvai) — a blanket}

\label{lesson:TA-C119-porvai}



\begin{quote}
Before the new one: say the Tamil for a towel; a piece, then the Tamil for a pillow.

Say the day words, \textbf{\ta{நாள்}} and the formal \textbf{\ta{தினம்}}, then the night words, \textbf{\ta{இரவு}}, \textbf{\ta{ராத்திரி}}, and \textbf{\ta{இருள்}}, darkness.
\end{quote}

\subsection*{You'll want to know: \ta{போர்வை}}
\textbf{\ta{போர்வை}} — \emph{pōrvai} — \textquotedblleft{}a blanket\textquotedblright{}.

\textbf{In the family, and next door}

\noindent
\begin{tabularx}{\linewidth}{@{}>{\raggedright\arraybackslash}X>{\raggedright\arraybackslash}X@{}}
\toprule
\textbf{} & \textbf{word} \\
\midrule
Kannada & \ka{ಕಂಬಳಿ} (\emph{kambaḷi}) \\
Telugu & \te{దుప్పటి} (\emph{duppaṭi}) \\
Malayalam & \ml{പുതപ്പ്} (\emph{putappŭ}) \\
Hindi (neighbour) & \dv{कंबल} (\emph{kambal}) \\
English & a blanket \\
\bottomrule
\end{tabularx}

\begin{grammarlens}[title={What you've built}]
One more word for this chapter.
\end{grammarlens}

\subsection*{Your turn}
\begin{itemize}
\raggedright
  \item \emph{Say it:} \emph{pōrvai}
  \item \emph{Say it:} \emph{pōrvai}, once more
  \item \emph{Say it:} say \emph{pōrvai}
  \item \emph{Recall:} say the Tamil for a towel; a piece, then the Tamil for a pillow, then say \emph{pōrvai} again
\end{itemize}

\begin{tcolorbox}[breakable,colback=teal!4,colframe=teal!35!black,title={Before you move on}]
What does \ta{போர்வை} mean? (\textquotedblleft{}A blanket\textquotedblright{}.) Say it once more.
\end{tcolorbox}
